\documentclass[12pt]{article}
\usepackage[utf8]{inputenc}
\usepackage[T1]{fontenc}
\usepackage{hyperref}
\title{Southern Uzbek Language Resources for K-12 Curriculum Development}
\author{Compiled by Claude Sonnet 5 (claude-sonnet-5) for Andr\'as Kornai}
\date{}
\begin{document}
\maketitle

\section{Introduction}
Southern Uzbek (ISO 639-3 uzs; Glottocode sout2699) is a Turkic language spoken mainly in Afghanistan, with speakers also reported in Pakistan. Estimates of speaker numbers vary by source. This report is based on three web searches and is short because little K-12 material turned up. Nothing below is a full textbook or curriculum that I could verify.

\section{Primary Educational Resources}
No downloadable primers or textbooks for Southern Uzbek were found in the search results. Reporting describes the policy context:
\begin{itemize}
\item Press coverage says the Afghan Ministry of Education produced Uzbek textbooks for grades 1--3 and planned materials up to grade 12. It also says Uzbek instruction began in government schools of Uzbek-majority provinces in 1996, was halted during the wars, and was revived in 2011. Some speakers report that textbooks were scarce and that Uzbek was later dropped from the curriculum. See \url{https://8am.media/eng/the-third-most-spoken-language-in-afghanistan-scattered-efforts-fall-short-in-enriching-uzbek-language/}.
\item In 2021, Uzbekistan's Foreign Ministry said that rumours of the Taliban removing Uzbek from the school curriculum were false: \url{https://m.kun.uz/en/news/2021/08/31/mfa-denies-rumors-that-taliban-excludes-the-study-of-uzbek-language-from-school-curriculum}. I found no information on the current (2025--2026) status.
\end{itemize}

\section{Linguistic and Reference Materials}
No dictionary or grammar specific to Southern Uzbek was found. General reference pages:
\begin{itemize}
\item Wikipedia, Southern Uzbek language: \url{https://en.wikipedia.org/wiki/Southern_Uzbek_language}
\item Glottolog entry: \url{https://glottolog.org/resource/languoid/id/sout2699}
\item Wikitongues language profile: \url{https://wikitongues.org/languages/uzs/}
\end{itemize}
Nothing on orthography conventions for Afghan Uzbek (for example Arabic script use) was found.

\section{Religious and Community Resources}
\begin{itemize}
\item Joshua Project profile of Southern Uzbek in Afghanistan (demographic information): \url{https://joshuaproject.net/people_groups/15756/af}
\item Global Recordings Network lists audio Bible stories and messages for Southern Uzbek (the script is not stated): \url{https://www.globalrecordings.net/en/language/uzs}
\item Historical note: the earliest known Uzbek Scripture translation, the four Gospels in Arabic script, was published in Leipzig in 1891 (\url{https://ibtrussia.org/node/871}). It is not specific to Afghanistan.
\end{itemize}

\section{Digital and Audio Resources}
The only audio resource found is the Global Recordings page above. Uzbek Bible apps found (Muqaddas Kitob, Khorazm Uzbek) use Latin and Cyrillic scripts and are not Afghan Southern Uzbek resources.

\section{Curriculum Development Context}
Evidence of Southern Uzbek classroom materials is thin and the policy status is uncertain. A curriculum project would probably need to develop materials from scratch, with local partners, and settle on a script (Arabic script is used in Afghanistan, but I found no source documenting the standard). Further targeted searches in Dari/Uzbek and in Afghan Ministry of Education archives are advised.

\end{document}